% ============================================================
%  BEAMER — CÁLCULO INTEGRAL: INTEGRALES ALGEBRAICAS Y CAMBIO DE VARIABLE
%  Clase de refuerzo · Ejercicios de la memoria de clase
%  Fuente: tex_files/clases_integrales_ruben_usm/memoria_clases.md
%  Autor: Ing. Gabriel Astudillo
%  Compilar:  latexmk -xelatex beamer_integrales_algebraicas.tex
%
%  Alcance: integrales indefinidas algebraicas, cambio de variable y
%  división de polinomios. Los ejercicios (§1–§4 de la memoria) se muestran
%  con el enunciado anclado arriba y espacio en blanco para resolver a mano.
% ============================================================

\documentclass[aspectratio=169]{beamer}
\usepackage{mathwizards-beamer}

\mwbeamer{Cálculo Integral}{Integrales algebraicas y cambio de variable}{Clase de refuerzo $\cdot$ Sustitución y división de polinomios}

% ── Comandos propios de esta presentación ────────────────────
\newcommand{\dx}{\,dx}
\newcommand{\du}{\,du}

\begin{document}

\begin{frame}[plain]
  \titlepage
\end{frame}

% ════════════════════════════════════════════════════════════
\section{Herramientas de la clase}
% ════════════════════════════════════════════════════════════

\begin{frame}{Ruta de la clase}
  \begin{mwpasos}[Herramientas que repasamos]
    \begin{itemize}
      \item Regla de potencia y linealidad de la integral.
      \item Propiedades de exponentes (negativos y fraccionarios).
      \item Triángulo de Pascal para expandir binomios.
      \item Método de cambio de variable (sustitución).
      \item División de polinomios antes de integrar.
    \end{itemize}
  \end{mwpasos}
  \vspace{2pt}
  \begin{mwextra}[Alcance]
    Integrales \textbf{algebraicas} indefinidas. \textbf{No} fracciones parciales,
    \textbf{no} integración por partes y \textbf{no} integrales trigonométricas.
  \end{mwextra}
\end{frame}

\begin{frame}{Regla de potencia y linealidad}
  \begin{block}{Fórmulas base}
    \[
      \int x^{n}\dx = \frac{x^{n+1}}{n+1}+C \quad (n\neq -1)
      \qquad
      \int k\,f(x)\dx = k\int f(x)\dx
    \]
    \[
      \int \bigl[f(x)\pm g(x)\bigr]\dx = \int f(x)\dx \pm \int g(x)\dx
      \qquad
      \int \frac{1}{x}\dx = \ln|x|+C
    \]
  \end{block}
  \vspace{2pt}
  \begin{mwextra}[Con exponentes fraccionarios]
    $\sqrt[m]{x^{n}} = x^{n/m}$: se integra igual y luego se puede volver a raíz.
    Ejemplo: $\displaystyle\int u^{1/2}\du = \frac{2}{3}u^{3/2}+C$.
  \end{mwextra}
\end{frame}

\begin{frame}{Propiedades de exponentes}
  \begin{block}{Del cuaderno (pizarra)}
    \[
      \left(\frac{a}{b}\right)^{-n}=\left(\frac{b}{a}\right)^{n}
      \qquad
      \frac{1}{a^{n}}=a^{-n}
      \qquad
      a^{-n}=\frac{1}{a^{n}}
    \]
    \[
      \frac{a\pm b\pm c\pm d}{n}=\frac{a}{n}\pm\frac{b}{n}\pm\frac{c}{n}\pm\frac{d}{n}
    \]
  \end{block}
  \vspace{2pt}
  \begin{mwpasos}[Al dividir potencias]
    $\dfrac{x^{m}}{x^{n}} = x^{m-n}$.
    Ejemplo: $\dfrac{x^{-2/3}}{x^{5/2}} = x^{-2/3-5/2} = x^{-19/6}$.
  \end{mwpasos}
\end{frame}

\begin{frame}{Triángulo de Pascal}
  \begin{block}{Coeficientes de $(a\pm b)^{n}$}
    \[
      \begin{array}{c@{\quad}l}
        n=3 & 1\;\;3\;\;3\;\;1\\[3pt]
        n=5 & 1\;\;5\;\;10\;\;10\;\;5\;\;1\\[3pt]
        n=6 & \boxed{1\;\;6\;\;15\;\;20\;\;15\;\;6\;\;1}
      \end{array}
    \]
  \end{block}
  \vspace{2pt}
  \begin{mwextra}[Signos alternados]
    En $(a-b)^{n}$ los signos se \textbf{alternan}. Ejemplo (fila $n=6$):
    \[
      (u-2)^{6}=u^{6}-12u^{5}+60u^{4}-160u^{3}+240u^{2}-192u+64
    \]
  \end{mwextra}
\end{frame}

\begin{frame}{Método de cambio de variable}
  \begin{block}{Se deshace la regla de la cadena}
    \[
      \int f\bigl(g(x)\bigr)\,g'(x)\dx = \int f(u)\du = F(u)+C,
      \qquad u=g(x),\quad du=g'(x)\dx
    \]
  \end{block}
  \vspace{2pt}
  \begin{mwpasos}[Procedimiento en 5 pasos]
    \begin{enumerate}
      \item Elegir la sustitución $u=g(x)$.
      \item Calcular $du=g'(x)\dx$ (despejar $\dx$ si conviene).
      \item Reescribir todo en $u$ (sin dejar ninguna $x$).
      \item Integrar respecto a $u$.
      \item Volver a $x$ y añadir $+C$.
    \end{enumerate}
  \end{mwpasos}
\end{frame}

\begin{frame}{División de polinomios}
  \begin{block}{Cuándo dividir: grado $N \geq$ grado $D$}
    \[
      \frac{N(x)}{D(x)} = Q(x) + \frac{R(x)}{D(x)},
      \qquad \deg R < \deg D
    \]
    \[
      \int \frac{N(x)}{D(x)}\dx = \int Q(x)\dx + \int \frac{R(x)}{D(x)}\dx
    \]
  \end{block}
  \vspace{2pt}
  \begin{alertblock}{Ojo con esto}
    $\int Q$ es polinómica (directa). La fracción $R/D$ suele resolverse por
    sustitución. \textbf{Fracciones parciales: todavía no.}
  \end{alertblock}
\end{frame}

\begin{frame}{Repertorio trigonométrico (solo referencia)}
  \begin{mwextra}[Para más adelante — sin ejercicios en esta clase]
    Con $a\neq 0$:
    \[
      \int\sin(ax)\dx=-\frac{\cos(ax)}{a}+C
      \qquad
      \int\cos(ax)\dx=\frac{\sin(ax)}{a}+C
    \]
    \[
      \int\sec^{2}(ax)\dx=\frac{\tan(ax)}{a}+C
      \qquad
      \int\csc^{2}(ax)\dx=-\frac{\cot(ax)}{a}+C
    \]
    \[
      \int\sec(ax)\tan(ax)\dx=\frac{\sec(ax)}{a}+C
      \qquad
      \int\csc(ax)\cot(ax)\dx=-\frac{\csc(ax)}{a}+C
    \]
  \end{mwextra}
\end{frame}

% ════════════════════════════════════════════════════════════
\section{Ejercicios de clase}
% ════════════════════════════════════════════════════════════

\begin{frame}{Ejercicio 1}
  \begin{mwejercicio}[Ejercicio 1 — Sustitución + Pascal \quad \dificil]
    \[
      \int \sqrt{(x^{2}+2)^{3}}\;\cdot\;x^{12}\dx
    \]
    \small
    \textbf{Sugerencia:} $u = x^{2}+2$, de donde $x^{2} = u-2$.
    Expande $(u-2)^{6}$ con la fila $n=6$ de Pascal.
  \end{mwejercicio}
\end{frame}

\begin{frame}{Ejercicio 2}
  \begin{mwejercicio}[Ejercicio 2 — Radicales anidados \quad \muyDificil]
    \[
      \int \sqrt{\,2-\sqrt{1+\sqrt{x}\,}\;}\;\dx
    \]
    \small
    \textbf{Sugerencia:} $u = 2-\sqrt{1+\sqrt{x}}$. Despeja encadenando:
    $\sqrt{1+\sqrt{x}} = 2-u$ y luego $\sqrt{x} = (2-u)^{2}-1$.
  \end{mwejercicio}
\end{frame}

\begin{frame}{Ejercicio 3}
  \begin{mwejercicio}[Ejercicio 3 — División de polinomios \quad \muyDificil]
    \[
      \int \frac{x^{10}+2x^{9}+x^{8}-x^{8}-2x^{5}+2x^{2}}{2x^{6}+4x^{3}+2}\dx
    \]
    \small
    \textbf{Sugerencia:} factoriza el denominador $2x^{6}+4x^{3}+2 = 2(x^{3}+1)^{2}$
    y divide los polinomios antes de integrar.
  \end{mwejercicio}
\end{frame}

\begin{frame}{Ejercicio 4}
  \begin{mwejercicio}[Ejercicio 4 — Exponentes fraccionarios \quad \muyDificil]
    \[
      \int \left(\frac{\tfrac{3}{2}x^{-2/3}-\tfrac{1}{2}x^{-3/4}}{\tfrac{1}{4}x^{5/2}}\right)^{2}\dx
    \]
    \small
    \textbf{Sugerencia:} divide término a término, resta los exponentes y eleva
    el binomio al cuadrado $(A-B)^{2}=A^{2}-2AB+B^{2}$.
  \end{mwejercicio}
\end{frame}

% ════════════════════════════════════════════════════════════
\section{Notas de la memoria}
% ════════════════════════════════════════════════════════════

\begin{frame}{Notas — Ejercicio 1}
  \begin{mwextra}[Obs.~1.1 — el factor $x$ del $\dx$]
    Con $\dx = \du/(2x)$, al multiplicar por $x^{12}$ debería quedar
    $\tfrac12 x^{11}\du$, no $\tfrac12 x^{12}\du$. Revisar el enunciado:
    ¿es $x^{12}$ o $x^{11}$? ¿y el exponente de la raíz ($3/2$ u otro)?
  \end{mwextra}
  \vspace{3pt}
  \begin{mwextra}[Obs.~1.2 — exponentes de la primitiva]
    Al integrar se suma $1$ (no $3$) a cada exponente: deben quedar
    $17/2,\,15/2,\,\dots,\,5/2$ (raíces de potencias $17,15,\dots,5$).
  \end{mwextra}
\end{frame}

\begin{frame}{Notas — Ejercicios 2 y 3}
  \begin{mwextra}[Obs.~2.1 — signo global y último coeficiente]
    La primitiva en $u$ lleva factor $-4$; y el último coeficiente es
    $6\cdot\tfrac23 = 4$ (no $\tfrac43$).
  \end{mwextra}
  \vspace{3pt}
  \begin{mwextra}[Obs.~3.1 — cociente y residuo]
    El cociente/residuo de la hoja ($x^{4}-x^{2}$ y $\tfrac{3x^{2}}{(x^{3}+1)^{2}}$)
    no sale directo del numerador escrito: rehacer la división larga. El método
    (sacar $\tfrac12$, dividir y sustituir $u=x^{3}+1$) está bien planteado.
  \end{mwextra}
\end{frame}

\begin{frame}{Notas — Ejercicio 4 (pizarra)}
  \begin{mwextra}[Obs.~4.1 — el enunciado de la pizarra]
    La primera línea con radicales está tachada/borrosa; usar la
    \textbf{versión de potencias}, que es autoconsistente con el desarrollo.
  \end{mwextra}
  \vspace{3pt}
  \begin{mwextra}[Obs.~4.2 — cierre esperado]
    El ejercicio queda inconcluso en la pizarra. Si se integra:
    \[
      -\tfrac{27}{4}x^{-16/3}+\tfrac{288}{65}x^{-65/12}-\tfrac{8}{11}x^{-11/2}+C
    \]
  \end{mwextra}
\end{frame}

% ════════════════════════════════════════════════════════════
\section{Para practicar}
% ════════════════════════════════════════════════════════════

\begin{frame}{Para practicar}
  \begin{mwpasos}[Guía de la clase — 30 ejercicios progresivos]
    \begin{itemize}
      \item \textbf{Bloque I:} cambio de variable directo (8 problemas).
      \item \textbf{Bloque II:} sustitución con despeje y potencias fraccionarias (10).
      \item \textbf{Bloque III:} división de polinomios previa (12).
      \item Incluye clave de respuestas.
    \end{itemize}
  \end{mwpasos}
  \vspace{2pt}
  \begin{mwextra}[Alcance de la guía]
    Algebraicas $+$ sustitución $+$ división de polinomios. Sin fracciones
    parciales, sin partes y sin trigonométricas.
  \end{mwextra}
\end{frame}

\end{document}
